\section{Evaluation Protocol}
\label{sec:protocol}

The protocol is designed so that a difference between arms can be attributed to the lever, not to the prompt, the model or the measurement. It fixes the sender, receiver, scaffold and decoder, measures cost on the prompt side and under serving conditions, and scores accuracy with a rule agreed before the results existed.

\subsection{Analysis Plan and Change Log}
The study was not recorded in a third-party study registry. Instead, the analysis plan, with its endpoints, comparators, statistics and decision rules, was committed to version control on 22 September 2026, before any run of the extended study. Every later change of design went into a numbered log, with its date and reason, before the data it affected existed; earlier text was never edited. Analyses run after seeing results are labelled exploratory. The change log has 29 numbered entries, some with lettered amendments; the last, a training-data scaling curve, was stopped before it produced any result. For the external test the order is stricter: the confirmatory plan was committed before the external database was downloaded, the system was then frozen by a manifest of file hashes and a release tag, and the hash of the test-window list was committed before any model output on that database existed. Commit timestamps are self-declared, so we release the complete version history, including the change log, the freeze manifest and every result file, in an archived snapshot, so that the order of plan, data and results can be checked commit by commit. We describe this as a time-stamped analysis plan held in version control. We say where a change of design altered an outcome, and where a rule fixed in advance gave an answer we do not read literally (Section~\ref{sec:results}).

\subsection{Data}
\textbf{Training and primary test.} The MIT-BIH Arrhythmia Database \cite{moody2001impact} was used with the inter-patient split of de~Chazal et al.\ \cite{dechazal2004automatic}: 22 records for training and 22 for testing, with no patient in both, after the four paced records were excluded. Of the training records, 19 were used for training and three held out entirely for validation, so that no tuning decision touched a test patient. Beats were extracted as 360-sample windows centred on the annotated R peak and standardised per window, and annotations were mapped to the five AAMI classes. The test set contains no unclassifiable beats, a property of the split.

\textbf{Chronological replay.} The abnormal-run count is state carried across beats of a record. Events were therefore produced by replaying each record in order from its start through the frozen perception agent, so that the run count and RR interval of every window are those a deployed system would have computed.

\textbf{Test windows.} A fixed sampler (seed 0) draws windows from the test records, stratified by reference tier and by the class of the window's most urgent beat, at most three per record per cell, so that urgent and priority windows are not dominated by one class. The primary set has 1,195 windows from 22 patients, of which 400 are drawn at natural prevalence to measure false alarms. Cells that cannot be filled are left short and the achieved counts are reported in Table~\ref{tab:windows}.

\begin{table}[t]
\caption{Test windows by number of events per prompt $N$ and reference tier (primary DS2 set, 22 records). Stratified windows are drawn per tier; natural-prevalence windows are drawn without stratification and measure false alarms at real tier frequencies. Cells shorter than their quota (for example priority at $N=50$) are those the test records could not fill.}
\label{tab:windows}
\centering
\footnotesize
\adjustbox{max width=\linewidth}{\begin{tabular}{rrrrrrr}
\toprule
$N$ & Routine & Priority & Urgent & Stratified & Natural prevalence & Total \\
\midrule
1 & 60 & 60 & 57 & 177 & --- & 177 \\
5 & 60 & 59 & 56 & 175 & 100 & 275 \\
10 & 60 & 53 & 56 & 169 & 100 & 269 \\
20 & 57 & 45 & 55 & 157 & 100 & 257 \\
50 & 40 & 37 & 40 & 117 & 100 & 217 \\
All &  &  &  & 795 & 400 & 1195 \\
\bottomrule
\end{tabular}
}
\end{table}

\textbf{External confirmatory data.} The St Petersburg INCART 12-lead arrhythmia database \cite{goldberger2000physionet} (75 half-hour recordings from 32 patients) was fixed as the confirmatory test before it was downloaded or inspected. Lead II is resampled from 257 to 360~Hz and converted with the same procedure, and no INCART data are used for training, tuning or prompt design. The adapter, the text prompts and the calibration biases are frozen and hashed in a manifest before the run. The results are reported in Section~\ref{sec:incart}.

\subsection{Comparators and Calibration}
The comparators are A-compact, A-full (cost only) and A-filtered (Table~\ref{tab:levers}). Reusing computation, the second lever, is not an arm: prefix caching is switched on and off for every arm in the serving measurements, to check the ranking of the other levers under it. A text model with every event listed is biased toward answering ``routine'' on long prompts, an effect known from few-shot prompting \cite{zhao2021calibrate}. To avoid a strawman we report each text arm under default decoding and with a bias added to the routine score at the tier step, chosen on the held-out validation windows only and then frozen. The calibrated arm is the fair comparison and is used for every inferential statement.

\subsection{Endpoints and Statistics}
\textbf{Accuracy.} The primary endpoint is balanced accuracy, the mean recall over the tiers present, so that the rare urgent tier carries the same weight as the common routine tier. Arms are compared on the same windows (paired). Because windows from one patient are correlated, confidence intervals use a patient-cluster bootstrap with 20,000 resamples and percentile 95\% limits. The primary question for the adapter against calibrated text is non-inferiority with a margin of $-0.05$, tested first, followed by superiority where non-inferiority holds, in a fixed order at $\alpha = 0.05$. Adapter results are the mean of three training seeds, and seed ranges are reported.

\textbf{False alarms.} The share of routine windows at natural prevalence that the arm escalates to priority or urgent.

\textbf{Cost.} Token counts of the prompt; time to first token (prompt processing) and time to the tier token (the decision), at batch size one with arm order rotated; context memory; and GPU energy. We report time to the decision, not time to the full answer, because answer length differs across arms while decode speed is identical (about 160~ms per token); the full answer would credit the interface with differences in the length of the justification. Cost medians use a paired bootstrap of the relative difference on 21 windows per $N$.

\textbf{Serving.} Throughput and energy per decision were measured with batches of one, four and eight, repeated three times, and with prefix caching on and off. The framework was capped at 90\% of the card's memory so that a batch that does not fit is recorded as such instead of being silently paged into system memory, which would measure paging and not the model. Energy is the GPU board power reported by the NVIDIA Management Library, sampled at 100~Hz in a background thread during each timed generation call (prompt processing and generation up to the tier token) and integrated over the call's wall-clock time with the trapezoid rule; energy per decision is the batch's energy divided by the batch size. The measurement covers the whole board, including its idle and static draw, and excludes the processor, host memory and the rest of the system; every call spanned at least 87 power samples. Each configuration was measured on every window chunk in each of the three repeats (63 measurements per cell at batch one, 15 at batch four and 6 at batch eight), and we report the median with the interquartile range. Serving was measured for one adapter (the first training seed), so the spread reflects measurement and content variation, not training variation.

\textbf{Recoverability.} For each event, a decoder is trained on the training records to predict label, tier, heart rate, RR interval and abnormal-run status from the virtual tokens of a given slot, and tested on the test records. The same decoder family is trained on the 32-dimensional input and on the 35-dimensional input including side inputs, so that the loss attributable to the channel is the gap between the two. Decoders are fitted per slot (slots 0, 9, 19 and 49). The decoder fixed in advance is a small multilayer perceptron; a linear decoder is fitted alongside it for comparison. Run status is scored as balanced accuracy for a run of three or more beats, because $R^2$ is ill-posed for the heavy-tailed run length on this split (a change recorded before any data from it were analysed).

\subsection{Implementation and Reproducibility}
All experiments ran on one workstation with an NVIDIA RTX 5070 (12~GB), Python 3.13 and HuggingFace \texttt{transformers} with \texttt{bitsandbytes}. Production serving engines do not accept continuous input embeddings, so absolute latencies are higher than an optimised engine would give; every arm shares this, so only relative comparisons are meaningful. Every run records its git commit, checkpoint hashes, seed and library versions. Adapter training checkpoints resume bit-identically, and one run was paused and resumed without effect on the result. The tables and figures are generated from the result files by scripts in the supplementary release. \todo{Add repository and archive DOI after the freeze.}
